\documentclass[10pt,a4paper]{amsart}
\usepackage{amsmath,amssymb,amsthm,graphicx}
% Retained minimal preamble for the Stacks Project chapter more-algebra.tex (More on Algebra).
% The chapter file pinned-original-source.tex carries only \input{preamble} on its first line, so no
% preamble text lives inside the chapter itself.  The Stacks Project uses one single shared preamble for
% every chapter; the snapshot stages it identically as sources/stacks-chow/preamble.tex and
% sources/stacks-etale-cohomology/preamble.tex.  The block below is copied mechanically, line for line,
% from sources/stacks-chow/preamble.tex lines 14-18 (the Xy-pic setup) and 154-217 (the native theorem
% declarations together with the mathematical macros).  Nothing was added, removed or redefined; the
% wrapper class and the amsmath/amssymb/amsthm/graphicx base packages are supplied by the assembler.
% Chapter: more-algebra.tex -- More on Algebra.
% For commutative diagrams we use Xy-pic
\usepackage[all]{xy}

% We use 2cell for 2-commutative diagrams.
\xyoption{2cell}

% Theorem environments.
%
\theoremstyle{plain}
\newtheorem{theorem}[subsection]{Theorem}
\newtheorem{proposition}[subsection]{Proposition}
\newtheorem{lemma}[subsection]{Lemma}

\theoremstyle{definition}
\newtheorem{definition}[subsection]{Definition}
\newtheorem{example}[subsection]{Example}
\newtheorem{exercise}[subsection]{Exercise}
\newtheorem{situation}[subsection]{Situation}

\theoremstyle{remark}
\newtheorem{remark}[subsection]{Remark}
\newtheorem{remarks}[subsection]{Remarks}

\numberwithin{equation}{subsection}

% Macros
%
\def\lim{\mathop{\mathrm{lim}}\nolimits}
\def\colim{\mathop{\mathrm{colim}}\nolimits}
\def\Spec{\mathop{\mathrm{Spec}}}
\def\Hom{\mathop{\mathrm{Hom}}\nolimits}
\def\Ext{\mathop{\mathrm{Ext}}\nolimits}
\def\SheafHom{\mathop{\mathcal{H}\!\mathit{om}}\nolimits}
\def\SheafExt{\mathop{\mathcal{E}\!\mathit{xt}}\nolimits}
\def\Sch{\mathit{Sch}}
\def\Mor{\mathop{\mathrm{Mor}}\nolimits}
\def\Ob{\mathop{\mathrm{Ob}}\nolimits}
\def\Sh{\mathop{\mathit{Sh}}\nolimits}
\def\NL{\mathop{N\!L}\nolimits}
\def\CH{\mathop{\mathrm{CH}}\nolimits}
\def\proetale{{pro\text{-}\acute{e}tale}}
\def\etale{{\acute{e}tale}}
\def\QCoh{\mathit{QCoh}}
\def\Ker{\mathop{\mathrm{Ker}}}
\def\Im{\mathop{\mathrm{Im}}}
\def\Coker{\mathop{\mathrm{Coker}}}
\def\Coim{\mathop{\mathrm{Coim}}}

% Boxtimes
%
\DeclareMathSymbol{\boxtimes}{\mathbin}{AMSa}{"02}

%
% Macros for moduli stacks/spaces
%
\def\QCohstack{\mathcal{QC}\!\mathit{oh}}
\def\Cohstack{\mathcal{C}\!\mathit{oh}}
\def\Spacesstack{\mathcal{S}\!\mathit{paces}}
\def\Quotfunctor{\mathrm{Quot}}
\def\Hilbfunctor{\mathrm{Hilb}}
\def\Curvesstack{\mathcal{C}\!\mathit{urves}}
\def\Polarizedstack{\mathcal{P}\!\mathit{olarized}}
\def\Complexesstack{\mathcal{C}\!\mathit{omplexes}}
% \Pic is the operator that assigns to X its picard group, usage \Pic(X)
% \Picardstack_{X/B} denotes the Picard stack of X over B
% \Picardfunctor_{X/B} denotes the Picard functor of X over B
\def\Pic{\mathop{\mathrm{Pic}}\nolimits}
\def\Picardstack{\mathcal{P}\!\mathit{ic}}
\def\Picardfunctor{\mathrm{Pic}}
\def\Deformationcategory{\mathcal{D}\!\mathit{ef}}

% --- appended 2026-10-12 (Chief): the project own definitions of the reference/slogan/history
% display environments, copied mechanically from the shared Stacks preamble lines 8-12
% (\usepackage{verbatim} plus the three \newenvironment lines, which the chapter preamble keeps
% verbatim from the project preamble). Nothing added, removed or redefined: these environments
% carry no mathematical content and the project itself typesets them as comments. Without these
% three lines a counted statement whose own text uses them fails with Environment slogan undefined.
\usepackage{verbatim}
\newenvironment{reference}{\comment}{\endcomment}
\newenvironment{slogan}{\comment}{\endcomment}
\newenvironment{history}{\comment}{\endcomment}
\usepackage{cleveref}
\crefname{originallemma}{original lemma}{original lemmas}
\Crefname{originallemma}{Original lemma}{Original lemmas}
\makeatletter
\newlabel{algebra-lemma-characterize-unramified}{{02FM}{source}{}{sourcepointer-1}{}}
\newlabel{algebra-lemma-characterize-unramified@cref}{{[originallemma][0][]02FM}{[1][0][]source}{}{}{}}
\makeatother
\crefname{originallemma}{original lemma}{original lemmas}
\Crefname{originallemma}{Original lemma}{Original lemmas}
\makeatletter
\newlabel{algebra-proposition-unramified-locally-standard}{{0395}{source}{}{sourcepointer-2}{}}
\newlabel{algebra-proposition-unramified-locally-standard@cref}{{[originallemma][0][]0395}{[1][0][]source}{}{}{}}
\makeatother
\crefname{originallemma}{original lemma}{original lemmas}
\Crefname{originallemma}{Original lemma}{Original lemmas}
\makeatletter
\newlabel{lemma-ideal-inverse-nonzero}{{31702}{source}{}{sourcepointer-3}{}}
\newlabel{lemma-ideal-inverse-nonzero@cref}{{[originallemma][0][]31702}{[1][0][]source}{}{}{}}
\makeatother
\crefname{originallemma}{original lemma}{original lemmas}
\Crefname{originallemma}{Original lemma}{Original lemmas}
\makeatletter
\newlabel{lemma-invariance-number-branches-smooth}{{31596}{source}{}{sourcepointer-4}{}}
\newlabel{lemma-invariance-number-branches-smooth@cref}{{[originallemma][0][]31596}{[1][0][]source}{}{}{}}
\makeatother
\crefname{originallemma}{original lemma}{original lemmas}
\Crefname{originallemma}{Original lemma}{Original lemmas}
\makeatletter
\newlabel{lemma-number-of-branches-1}{{31566}{source}{}{sourcepointer-5}{}}
\newlabel{lemma-number-of-branches-1@cref}{{[originallemma][0][]31566}{[1][0][]source}{}{}{}}
\makeatother
\crefname{originallemma}{original lemma}{original lemmas}
\Crefname{originallemma}{Original lemma}{Original lemmas}
\makeatletter
\newlabel{lemma-reduced-goes-up}{{10533}{source}{}{sourcepointer-6}{}}
\newlabel{lemma-reduced-goes-up@cref}{{[originallemma][0][]10533}{[1][0][]source}{}{}{}}
\makeatother
\title{An essentially finite type injective local homomorphism out of a geometrically unibranch local domain with separable residue extension is the localization of an etale map}
\author{The Stacks Project Authors}
\date{stacks-project:more-algebra.tex:local-snapshot-20261012; GNU FDL 1.2 or later}
\begin{document}
\maketitle
\setcounter{subsection}{1}
\noindent\textit{Editorial note.} All delivered mathematical text is the author's own, copied from the pinned Stacks Project chapter file more-algebra.tex; no statement or proof step was rewritten, reordered or removed.\par\medskip
\noindent\textit{Editorial note.} The chapter file pinned-original-source.tex carries only \textbackslash{}input\{preamble\} on its first line, so no preamble text lives inside the chapter. The retained file sources/stacks-more-algebra/retained-preamble.tex is a mechanical, line-for-line copy of the project's single shared preamble as staged in the snapshot, together with the project's own definitions of the reference/slogan/history display environments, which the project itself typesets as comments; nothing was added, removed or redefined by this card and no package was added.\par\medskip
\noindent\textit{Editorial note.} The source URL is the source-level Stacks Project URL, as required when more than one card is produced from one source; the per-claim tag is carried by claim\_key.\par\medskip
\noindent\textit{Editorial note.} The counted statement span includes the author's own \textbackslash{}begin\{reference\} block recording that the result generalizes a cited theorem; it is delivered verbatim and produces no unresolved citation. The author writes 'see by Lemmas ...' inside the counted proof span, reproduced unchanged. The six references inside the counted proof span are recorded as non-load-bearing navigation pointers.\par\medskip
\noindent\textit{Allowed standard prerequisites (editorial).}\begin{itemize}
\item The characterisation of unramified ring maps and the structure of an unramified map as a quotient of a standard etale map. invoked at proof lines 31834-31839. Produces the standard etale cover C' and the surjection onto C.
\item Geometric unibranchness and the invariance of the number of branches under smooth maps. used at proof lines 31847-31849. Is what makes the standard etale cover a domain, the key step of the argument.
\item Nonvanishing of the contraction of a nonzero ideal for a flat localized etale map of domains. invoked at proof line 31851. Produces the contradiction that forces injectivity of B' -> B.
\end{itemize}\medskip
\ifdefined\hypertarget\hypertarget{sourcepointer-1}{}\fi
\noindent\textit{Source reference.} Original lemma 02FM: focused-seeds/AGP/sources/stacks-snapshot-20261011/tags.txt, source lines 3197--3197. Cross-chapter navigation pointer into the Stacks Project; the standard fact it supplies is applied inside the counted proof span and the argument resting on it stays there.\par\medskip
\ifdefined\hypertarget\hypertarget{sourcepointer-2}{}\fi
\noindent\textit{Source reference.} Original lemma 0395: focused-seeds/AGP/sources/stacks-snapshot-20261011/tags.txt, source lines 4260--4260. Cross-chapter navigation pointer into the Stacks Project; the standard fact it supplies is applied inside the counted proof span and the argument resting on it stays there.\par\medskip
\ifdefined\hypertarget\hypertarget{sourcepointer-3}{}\fi
\noindent\textit{Source reference.} Original lemma 31702: sources/stacks-more-algebra/pinned-original-source.tex, source lines 31702--31702. Same-chapter navigation pointer; the statement it names is standard for the section and the argument that consumes it stays inside the counted proof span.\par\medskip
\ifdefined\hypertarget\hypertarget{sourcepointer-4}{}\fi
\noindent\textit{Source reference.} Original lemma 31596: sources/stacks-more-algebra/pinned-original-source.tex, source lines 31596--31596. Same-chapter navigation pointer; the statement it names is standard for the section and the argument that consumes it stays inside the counted proof span.\par\medskip
\ifdefined\hypertarget\hypertarget{sourcepointer-5}{}\fi
\noindent\textit{Source reference.} Original lemma 31566: sources/stacks-more-algebra/pinned-original-source.tex, source lines 31566--31566. Same-chapter navigation pointer; the statement it names is standard for the section and the argument that consumes it stays inside the counted proof span.\par\medskip
\ifdefined\hypertarget\hypertarget{sourcepointer-6}{}\fi
\noindent\textit{Source reference.} Original lemma 10533: sources/stacks-more-algebra/pinned-original-source.tex, source lines 10533--10533. Same-chapter navigation pointer; the statement it names is standard for the section and the argument that consumes it stays inside the counted proof span.\par\medskip
\begin{lemma}
\label{lemma-unramified-extension-unibranch-domain-is-etale}
\begin{reference}
Generalization of \cite[Expose I, Theorem 9.5 part (ii)]{SGA1}
\end{reference}
Let $(A, \mathfrak m)$ be a geometrically unibranch local domain. Let
$A \to B$ be an injective local homomorphism of local rings,
which is essentially of finite type. If $\mathfrak m B$
is the maximal ideal of $B$ and the induced extension of
residue fields is separable, then $A \to B$ is the localization
of an \'etale ring map.
\end{lemma}

\begin{proof}
We may write $B = C_\mathfrak q$ where $A \to C$ is a finite
type ring map and $\mathfrak q \subset C$ is a prime ideal
lying over $\mathfrak m$. By
Algebra, Lemma \ref{algebra-lemma-characterize-unramified}
the ring map $A \to C$ is unramified at $\mathfrak q$.
By Algebra, Proposition \ref{algebra-proposition-unramified-locally-standard}
after replacing $C$ by a principal localization,
we can find a standard \'etale ring map $A \to C'$
and a surjection $C' \to C$. Denote $\mathfrak q' \subset C'$
the inverse image of $\mathfrak q$ and set $B' = C'_{\mathfrak q'}$.
Then $B' \to B$ is surjective.
It suffices to show that $B' \to B$ is also injective.

\medskip\noindent
Since $A$ is a domain, the rings $C'$ and $B'$ are reduced, see
Algebra, Lemma \ref{lemma-reduced-goes-up}.
Since $A$ is geometrically unibranch, the ring $B'$ is a domain, see
by Lemmas \ref{lemma-invariance-number-branches-smooth} and
\ref{lemma-number-of-branches-1}.
If $B' \to B$ is not injective, then
$A \cap \Ker(B' \to B)$ is nonzero by Lemma \ref{lemma-ideal-inverse-nonzero}
which contradicts the assumption that $A \to B$ is injective.
\end{proof}

\end{document}
